% =====================================================================
% Section: Discussion  (outline.md §4)
% Draft v1 (2026-06-27). Discusses the applicability boundary of content routing,
%   the risk of over-strong routing, prototype-count trade-off, and limitations.
% Consistent with the honest framing of the experiments:
%   - DPRNet matches/slightly surpasses CATANet at comparable compute
%     (3_experiments_comparison.tex, 3_experiments_efficiency.tex).
%   - Ablation is now a 250k, 3-seed study (mean+/-std): single-switch
%     attribution + C1 direct EMA comparison + C4 variance reduction are supported;
%     margins are still small (3_experiments_ablation.tex,
%     method-experiment-traceability.md).
% No new numeric claims beyond those already established in the body.
% =====================================================================
\section{Discussion}
\label{sec:discussion}

\paragraph{Content routing vs.\ spatial locality.}
The benchmark results in Section~\ref{sec:main_comparison} indicate where content-aware
routing helps most. On Urban100 and Manga109---images dominated by repetitive structures
and long-range self-similarity---DPRNet and CATANet, the two routing methods, separate
more noticeably from the window- and distillation-based baselines, whereas on the more
stochastic textures of BSD100 the margin between all top-tier methods narrows. This is
consistent with the premise of content routing and with earlier SR work that
exploits internal self-similarity~\cite{glasner2009single} or cross-scale
non-local correspondences~\cite{mei2020csnln}: grouping tokens by semantics
is most useful when semantically related content is spatially scattered. It is less
useful when local statistics already capture the signal. DPR does not discard spatial
locality; it complements it, since each TAB is paired with a local self-attention unit
(Section~\ref{sec:method_overview}); the routing module supplies the long-range, content-based
interaction that a fixed spatial window cannot.

\paragraph{Small margins and evidence strength.}
The numerical margins over CATANet are deliberately framed as small. The value of DPR is
therefore not a claim of large PSNR gain over its direct predecessor, but a combination of
three evidence streams: a consistent top-tier position across five benchmarks and three
scales (Section~\ref{sec:main_comparison}), a preserved lightweight profile relative to
attention-heavy Transformer baselines (Section~\ref{sec:efficiency}), and explicit routing
diagnostics showing meaningful confidence scores and more balanced prototype usage
(Section~\ref{sec:routing_analysis}). This framing matters because modern lightweight
SR methods are often separated by only a few hundredths of a dB under standard
PSNR/SSIM evaluation; without mechanism-level analysis, such small differences would be
weak evidence on their own.

\paragraph{The risk of over-strong routing.}
A confidence-aware router can, in principle, sharpen its assignments until tokens collapse
onto a few prototypes, which would defeat the purpose of routing and destabilize training.
Two components of DPR are explicitly designed against this failure mode: the routing
temperature is clamped (Eq.~\eqref{eq:scores}) and the usage-balancing regularizer
(Eq.~\eqref{eq:balance}) penalizes peaky usage distributions. The diagnostics in
Section~\ref{sec:routing_analysis} show both mechanisms operating as intended---the learned
temperature settles in a moderate, unsaturated range (block-mean $6.06$, below the clamp
of $10$) and the balancing regularizer raises the prototype-usage entropy in all eight
blocks ($0.623\!\rightarrow\!0.686$) while increasing the number of active prototypes.
Thus, the benefits of confidence-aware routing are obtained without letting the
router become degenerate.

\paragraph{Prototype count and the cost trade-off.}
DPR uses a per-block prototype budget $M=[16,32,64,128,16,32,64,128]$, growing within
each stage. More prototypes give the router a finer content vocabulary but enlarge the
query bank and the matching cost, while too few prototypes force unrelated tokens to
share a slot. We keep the prototype schedule aligned with the CATANet backbone so that
the comparison isolates the routing mechanism rather than a capacity change; this is also
why DPRNet carries a modestly higher parameter count and Multi-Adds than CATANet
($660$K/$52.14$\,G vs.\ $535$K/$46.8$\,G at $\times4$ under a common $256{\times}256$ input),
as the added cost lives in the query bank and projections of DPR. Even so, DPRNet stays
within the lightweight bracket and below the Transformer baselines SwinIR-light
($60.3$\,G) and SRFormer-light ($56.5$\,G) while reaching higher Urban100 PSNR
(Section~\ref{sec:efficiency}).
A systematic sweep of the prototype schedule is a natural extension but is orthogonal
to the routing formulation studied here.

\paragraph{Limitations.}
Several limitations bound the present study and frame our claims conservatively. The
ablations (Section~\ref{sec:ablation}) train every variant over three seeds, which
supports single-switch attribution and lets us report both the mean effect and its
seed-to-seed variance; the margins over the strongest configuration nonetheless remain small,
so we frame the design elements as mutually compatible and at least quality-preserving
rather than as sources of large individual gains. The direct comparison against a re-implemented EMA
center branch (C1) is conducted only at $\times4$ on Set5 and Urban100, the setting in which
content routing is most stressed, rather than across all five benchmarks. The efficiency
comparison is intentionally limited to common-protocol MACs and our measured DPRNet runtime:
we do not claim cross-method latency because the competing models were not remeasured on the
same hardware and implementation stack (Section~\ref{sec:efficiency}). The models are trained
on DIV2K only; larger training sets (e.g.\ DF2K) and real-world or arbitrary-scale
degradations are not explored here and may further shift the trade-off. Extending
confidence-aware routing to such settings---for instance alongside one-step diffusion
distillation~\cite{he2026tadsr} or dynamic autoregressive modeling~\cite{zheng2026dvar}
for real-world SR---is a promising direction that we leave to future work.
